
\documentclass{article}
\usepackage{colortbl}
\usepackage{makecell}
\usepackage{multirow}
\usepackage{supertabular}

\begin{document}

\newcounter{utterance}

\centering \large Interaction Transcript for game `wordle', experiment `medium\_frequency\_words\_no\_clue\_no\_critic', episode 9 with deepseek{-}v4{-}pro.
\vspace{24pt}

{ \footnotesize  \setcounter{utterance}{1}
\setlength{\tabcolsep}{0pt}
\begin{supertabular}{c@{$\;$}|p{.15\linewidth}@{}p{.15\linewidth}p{.15\linewidth}p{.15\linewidth}p{.15\linewidth}p{.15\linewidth}}
    \# & \multicolumn{2}{c}{Player} && \multicolumn{2}{c}{Game Master} \\
    \hline

    \theutterance \stepcounter{utterance}  
    & & & \multicolumn{4}{p{0.6\linewidth}}{
        \cellcolor[rgb]{0.9,0.9,0.9}{
            \makecell[{{p{\linewidth}}}]{
                \texttt{\tiny{[P1$\langle$GM]}}
                \texttt{Вы — языковой волшебник, который любит угадывать слова, используя заданные правила.} \\
\\ 
\texttt{Добро пожаловать в Wordle! У вас есть шесть попыток, чтобы угадать целевое слово — допустимое слово из пяти букв. Пожалуйста, используйте теги "объяснение:" и "догадка:", чтобы дать краткое объяснение для каждой попытки.} \\
\\ 
\texttt{Например, если ваша догадка — "слово", ваш ответ должен быть} \\
\texttt{объяснение: это распространённое слово из пяти букв, и я начинаю угадывать с этого слова.} \\
\texttt{догадка: слово} \\
\\ 
\texttt{После каждой попытки ваш ответ будет проверен, и вы получите обратную связь, показывающую, какие буквы верны (зелёный), какие буквы верны, но стоят не на своём месте (жёлтый), и какие буквы неверны (красный). Эта обратная связь может быть полезна для определения того, какие буквы включить или исключить в следующей попытке.} \\
\\ 
\texttt{Например, обратная связь для "слово" может быть:} \\
\texttt{оценка\_догадки: с<жёлтый> л<жёлтый> о<зелёный> в<жёлтый> о<красный>} \\
\\ 
\texttt{В explanation должны быть детали о том, как оценка\_догадки используется, чтобы прийти к новой догадке.} \\
\\ 
\texttt{Не генерируйте никакого другого текста.} \\
\\ 
\texttt{Начнём с вашей первой догадки.} \\
            }
        }
    }
     \\ \\

    \theutterance \stepcounter{utterance}  
    & \multicolumn{4}{p{0.6\linewidth}}{
        \cellcolor[rgb]{0.561,0.737,0.561}{
            \makecell[{{p{\linewidth}}}]{
                \texttt{\tiny{[P1$\rangle$GM]}}
                \texttt{объяснение: я начинаю с распространённого слова из пяти букв, чтобы охватить несколько часто используемых гласных и согласных.  } \\
\texttt{догадка: слово} \\
            }
        }
    }
    & & \\ \\

    \theutterance \stepcounter{utterance}  
    & & & \multicolumn{2}{p{0.3\linewidth}}{
        \cellcolor[rgb]{0.9,0.9,0.9}{
            \makecell[{{p{\linewidth}}}]{
                \texttt{\tiny{[GM$|$GM]}}
                \texttt{attempts = 1} \\
\texttt{target\_word = отдых} \\
\texttt{guess = слово} \\
\texttt{guess\_feedback = с<red> л<red> о<yellow> в<red> о<red>} \\
            }
        }
    }
    & & \\ \\

    \theutterance \stepcounter{utterance}  
    & & & \multicolumn{4}{p{0.6\linewidth}}{
        \cellcolor[rgb]{0.9,0.9,0.9}{
            \makecell[{{p{\linewidth}}}]{
                \texttt{\tiny{[P1$\langle$GM]}}
                \texttt{оценка\_догадки: с<красный> л<красный> о<жёлтый> в<красный> о<красный>} \\
\\ 
\texttt{Дайте ответ только в этом формате.} \\
\texttt{объяснение: подробности} \\
\texttt{догадка: слово} \\
            }
        }
    }
     \\ \\

    \theutterance \stepcounter{utterance}  
    & \multicolumn{4}{p{0.6\linewidth}}{
        \cellcolor[rgb]{0.561,0.737,0.561}{
            \makecell[{{p{\linewidth}}}]{
                \texttt{\tiny{[P1$\rangle$GM]}}
                \texttt{объяснение: буква "о" есть в слове, но не на третьем и не на пятом месте. Остальные буквы "с", "л", "в" исключены. Попробую новое слово с "о" на другой позиции и другими распространёнными буквами.  } \\
\texttt{догадка: тропа} \\
            }
        }
    }
    & & \\ \\

    \theutterance \stepcounter{utterance}  
    & & & \multicolumn{2}{p{0.3\linewidth}}{
        \cellcolor[rgb]{0.9,0.9,0.9}{
            \makecell[{{p{\linewidth}}}]{
                \texttt{\tiny{[GM$|$GM]}}
                \texttt{attempts = 2} \\
\texttt{target\_word = отдых} \\
\texttt{guess = тропа} \\
\texttt{guess\_feedback = т<yellow> р<red> о<yellow> п<red> а<red>} \\
            }
        }
    }
    & & \\ \\

    \theutterance \stepcounter{utterance}  
    & & & \multicolumn{4}{p{0.6\linewidth}}{
        \cellcolor[rgb]{0.9,0.9,0.9}{
            \makecell[{{p{\linewidth}}}]{
                \texttt{\tiny{[P1$\langle$GM]}}
                \texttt{оценка\_догадки: т<жёлтый> р<красный> о<жёлтый> п<красный> а<красный>} \\
\\ 
\texttt{Дайте ответ только в этом формате.} \\
\texttt{объяснение: подробности} \\
\texttt{догадка: слово} \\
            }
        }
    }
     \\ \\

    \theutterance \stepcounter{utterance}  
    & \multicolumn{4}{p{0.6\linewidth}}{
        \cellcolor[rgb]{0.561,0.737,0.561}{
            \makecell[{{p{\linewidth}}}]{
                \texttt{\tiny{[P1$\rangle$GM]}}
                \texttt{объяснение: буквы "т" и "о" есть в слове, но не на первой и третьей позициях соответственно, "р", "п", "а" исключены. Попробую разместить "т" и "о" на других местах и добавить новые буквы.  } \\
\texttt{догадка: поток} \\
            }
        }
    }
    & & \\ \\

    \theutterance \stepcounter{utterance}  
    & & & \multicolumn{2}{p{0.3\linewidth}}{
        \cellcolor[rgb]{0.9,0.9,0.9}{
            \makecell[{{p{\linewidth}}}]{
                \texttt{\tiny{[GM$|$GM]}}
                \texttt{attempts = 3} \\
\texttt{target\_word = отдых} \\
\texttt{guess = поток} \\
\texttt{guess\_feedback = п<red> о<yellow> т<yellow> о<red> к<red>} \\
            }
        }
    }
    & & \\ \\

    \theutterance \stepcounter{utterance}  
    & & & \multicolumn{4}{p{0.6\linewidth}}{
        \cellcolor[rgb]{0.9,0.9,0.9}{
            \makecell[{{p{\linewidth}}}]{
                \texttt{\tiny{[P1$\langle$GM]}}
                \texttt{оценка\_догадки: п<красный> о<жёлтый> т<жёлтый> о<красный> к<красный>} \\
\\ 
\texttt{Дайте ответ только в этом формате.} \\
\texttt{объяснение: подробности} \\
\texttt{догадка: слово} \\
            }
        }
    }
     \\ \\

    \theutterance \stepcounter{utterance}  
    & \multicolumn{4}{p{0.6\linewidth}}{
        \cellcolor[rgb]{0.561,0.737,0.561}{
            \makecell[{{p{\linewidth}}}]{
                \texttt{\tiny{[P1$\rangle$GM]}}
                \texttt{объяснение: буквы "о" и "т" есть в слове, но "о" не на втором и четвёртом месте, а "т" не на третьем. Кроме того, в слове только одна "о", так как вторая была красной. "п" и "к" исключены. Попробую слово, где "т" стоит на последнем месте, а "о" на первом.  } \\
\texttt{догадка: отчёт} \\
            }
        }
    }
    & & \\ \\

    \theutterance \stepcounter{utterance}  
    & & & \multicolumn{2}{p{0.3\linewidth}}{
        \cellcolor[rgb]{0.9,0.9,0.9}{
            \makecell[{{p{\linewidth}}}]{
                \texttt{\tiny{[GM$|$GM]}}
                \texttt{attempts = 4} \\
\texttt{target\_word = отдых} \\
\texttt{guess = отчёт} \\
\texttt{guess\_feedback = о<green> т<green> ч<red> ё<red> т<red>} \\
            }
        }
    }
    & & \\ \\

    \theutterance \stepcounter{utterance}  
    & & & \multicolumn{4}{p{0.6\linewidth}}{
        \cellcolor[rgb]{0.9,0.9,0.9}{
            \makecell[{{p{\linewidth}}}]{
                \texttt{\tiny{[P1$\langle$GM]}}
                \texttt{оценка\_догадки: о<зелёный> т<зелёный> ч<красный> ё<красный> т<красный>} \\
\\ 
\texttt{Дайте ответ только в этом формате.} \\
\texttt{объяснение: подробности} \\
\texttt{догадка: слово} \\
            }
        }
    }
     \\ \\

    \theutterance \stepcounter{utterance}  
    & \multicolumn{4}{p{0.6\linewidth}}{
        \cellcolor[rgb]{0.561,0.737,0.561}{
            \makecell[{{p{\linewidth}}}]{
                \texttt{\tiny{[P1$\rangle$GM]}}
                \texttt{объяснение: слово начинается на "от", "ч", "ё" и вторая "т" исключены. Нужно подобрать слово из пяти букв, начинающееся на "от", с новыми буквами вместо исключённых.  } \\
\texttt{догадка: отзыв} \\
            }
        }
    }
    & & \\ \\

    \theutterance \stepcounter{utterance}  
    & & & \multicolumn{2}{p{0.3\linewidth}}{
        \cellcolor[rgb]{0.9,0.9,0.9}{
            \makecell[{{p{\linewidth}}}]{
                \texttt{\tiny{[GM$|$GM]}}
                \texttt{attempts = 5} \\
\texttt{target\_word = отдых} \\
\texttt{guess = отзыв} \\
\texttt{guess\_feedback = о<green> т<green> з<red> ы<green> в<red>} \\
            }
        }
    }
    & & \\ \\

    \theutterance \stepcounter{utterance}  
    & & & \multicolumn{4}{p{0.6\linewidth}}{
        \cellcolor[rgb]{0.9,0.9,0.9}{
            \makecell[{{p{\linewidth}}}]{
                \texttt{\tiny{[P1$\langle$GM]}}
                \texttt{оценка\_догадки: о<зелёный> т<зелёный> з<красный> ы<зелёный> в<красный>} \\
\\ 
\texttt{Дайте ответ только в этом формате.} \\
\texttt{объяснение: подробности} \\
\texttt{догадка: слово} \\
            }
        }
    }
     \\ \\

    \theutterance \stepcounter{utterance}  
    & \multicolumn{4}{p{0.6\linewidth}}{
        \cellcolor[rgb]{0.561,0.737,0.561}{
            \makecell[{{p{\linewidth}}}]{
                \texttt{\tiny{[P1$\rangle$GM]}}
                \texttt{объяснение: слово состоит из букв "о", "т", "ы" на известных позициях: первая, вторая и четвёртая. "з" и "в" исключены. Осталось найти букву для третьего места, это должно быть допустимое слово. Подходит "отым" или "отык", но "отым" не является словом, а "отык" тоже маловероятно. Однако, вероятно, это слово "отым", которое в русском языке есть как форма глагола, но не уверен. Попробую другое: возможно, "отёк", но "ё" была красной ранее, так что нет. "отёк" не подходит из{-}за "ё". Значит, ищем слова с корнем "от...ы...". Возможно, "отит", но "т" в конце была красной, так что нет. Может быть, "отым" — форма "отнимать"? Это устаревшее или диалектное. Рискну.  } \\
\texttt{догадка: отым} \\
            }
        }
    }
    & & \\ \\

    \theutterance \stepcounter{utterance}  
    & & & \multicolumn{2}{p{0.3\linewidth}}{
        \cellcolor[rgb]{0.9,0.9,0.9}{
            \makecell[{{p{\linewidth}}}]{
                \texttt{\tiny{[GM$|$GM]}}
                \texttt{Error: The length of the guessed word is not 5.} \\
            }
        }
    }
    & & \\ \\

    \theutterance \stepcounter{utterance}  
    & & & \multicolumn{2}{p{0.3\linewidth}}{
        \cellcolor[rgb]{0.9,0.9,0.9}{
            \makecell[{{p{\linewidth}}}]{
                \texttt{\tiny{[GM$|$GM]}}
                \texttt{game\_result = ABORT} \\
            }
        }
    }
    & & \\ \\

\end{supertabular}
}

\end{document}
